% ============================================================
%  Basado en la plantilla de la Práctica 2 (mismos colores/estilos)
% ============================================================
\documentclass[11pt,a4paper]{article}

% ---------- Idioma y codificación ----------
\usepackage[utf8]{inputenc}
\usepackage[T1]{fontenc}
\usepackage[spanish,es-noshorthands]{babel}

% ---------- Tipografía y layout ----------
\usepackage{lmodern}
\usepackage[sfdefault,light]{roboto}
\usepackage{geometry}
\geometry{a4paper, top=3.1cm, bottom=2.8cm, left=2.4cm, right=2.4cm, headheight=1.6cm}
\usepackage{ragged2e}
\usepackage{enumitem}
\usepackage{graphicx}
\usepackage{float}
\usepackage{booktabs}
\usepackage{array}
\usepackage[most]{tcolorbox}
\usepackage{tikz}
\usetikzlibrary{shadows,calc}
\usepackage{titlesec}
\usepackage{fancyhdr}
\usepackage{fontawesome5}
\usepackage[colorlinks=true,breaklinks=true]{hyperref}
\usepackage{caption}
\usepackage{microtype}
\usepackage{xurl}

% ---------- Paleta de colores (tema "coolnight" de tu terminal Kitty) ----------
\definecolor{cNavy}{HTML}{011423}
\definecolor{cTealDeep}{HTML}{064B5A}
\definecolor{cTeal}{HTML}{097289}
\definecolor{cBlue}{HTML}{0FC5ED}
\definecolor{cCyan}{HTML}{0FC5ED}
\definecolor{cMint}{HTML}{44FFB1}
\definecolor{cPurple}{HTML}{A277FF}
\definecolor{cRed}{HTML}{E52E2E}
\definecolor{cBg}{HTML}{EAF6F8}
\definecolor{cText}{HTML}{01111E}

\hypersetup{
  linkcolor=cTealDeep,
  citecolor=cTealDeep,
  urlcolor=cBlue,
  pdftitle={Práctica 3 -- Pipeline de Clasificación con Weka},
  pdfauthor={Miguel Deza}
}

\color{cText}
\emergencystretch=2.5em

% ---------- Encabezado / pie de página ----------
\pagestyle{fancy}
\fancyhf{}
\renewcommand{\headrulewidth}{0pt}
\renewcommand{\footrulewidth}{0pt}
\fancyhead[L]{\small\color{cTealDeep}\textbf{\faGraduationCap\ Maestría en Ciencia de Datos}}
\fancyhead[R]{\small\color{cTeal}Práctica 3 -- Weka}
\fancyfoot[C]{%
  \begin{tikzpicture}[remember picture]
    \node[circle, fill=cTealDeep, minimum size=1.7em, inner sep=0pt] (pg) {\color{white}\small\thepage};
  \end{tikzpicture}%
}
\renewcommand{\headrule}{{\color{cCyan}\hrule height 1.4pt}\vspace{1pt}{\color{cBlue}\hrule height 0.5pt}}

% ---------- Estilo de secciones ----------
\titleformat{\section}
  {\normalfont\Large\bfseries\color{cTealDeep}}
  {\tikz[baseline=-0.6ex]{\node[fill=cCyan, text=white, rounded corners=1.5pt, inner sep=3pt, font=\bfseries]{\thesection};}}
  {0.6em}{}
  [{\vspace{-0.3em}\color{cCyan}\titlerule[1.2pt]\vspace{0.3em}}]

\titleformat{\subsection}
  {\normalfont\large\bfseries\color{cTeal}}
  {\color{cTeal}\thesubsection}{0.6em}{}

\titlespacing*{\section}{0pt}{1.4em}{0.8em}
\titlespacing*{\subsection}{0pt}{1.1em}{0.5em}

% ---------- Cajas reutilizables ----------
\newtcolorbox{objetivobox}{
  colback=cBg, colframe=cCyan, boxrule=0.9pt, arc=2mm,
  left=8pt, right=8pt, top=6pt, bottom=6pt,
  title={\faBullseye\ \ Objetivo}, coltitle=white,
  colbacktitle=cTealDeep, fonttitle=\bfseries,
  before upper={\raggedright\color{cText}}
}

\newtcolorbox{incidenciabox}{
  colback=white, colframe=cRed, boxrule=1pt, arc=1mm,
  left=8pt, right=8pt, top=6pt, bottom=6pt,
  title={\faExclamationTriangle\ \ Hallazgo técnico}, coltitle=white,
  colbacktitle=cRed, fonttitle=\bfseries,
  before upper={\raggedright\color{cText}}
}

\newtcolorbox{evidenciabox}{
  colback=cBg, colframe=cMint, boxrule=0.9pt, arc=2mm,
  left=8pt, right=8pt, top=6pt, bottom=6pt,
  title={\faFileArchive\ \ Evidencia de ejecución}, coltitle=white,
  colbacktitle=cTealDeep, fonttitle=\bfseries,
  before upper={\raggedright\color{cText}}
}

% ---------- Figura sin marco ----------
\newcommand{\figuraMarcada}[3]{%
  \begin{figure}[H]
    \centering
    \includegraphics[width=#2\linewidth]{#1}
    \caption{#3}
  \end{figure}
}

\captionsetup{labelfont={bf,color=cTealDeep}, font=small, textfont=it}

% ============================================================
\begin{document}

% -------------------- PORTADA --------------------
\begin{titlepage}
\newgeometry{margin=0cm}
\begin{tikzpicture}[remember picture, overlay]
  \fill[cNavy] (current page.south west) rectangle (current page.north east);

  \fill[cPurple, opacity=0.28] (current page.north east) ++(-1,1) circle (3.2);
  \fill[cCyan, opacity=0.22] (current page.north east) ++(-2.6,0.4) circle (2.1);
  \fill[cTealDeep, opacity=0.55] (current page.south west) ++(1,-1) circle (4.0);
  \fill[cMint, opacity=0.16] (current page.south west) ++(2.4,-0.2) circle (2.3);

  % ---- Icono insignia ----
  \node[circle, fill=cTealDeep, draw=cCyan, line width=1.3pt, minimum size=2.5cm] (logo) at ($(current page.north)+(0,-2.9)$) {};
  \node at (logo) {\color{cMint}\fontsize{30}{30}\selectfont\faBrain};

  \node[anchor=north, text width=14cm, align=center] at ($(current page.north)+(0,-4.7)$) {
    {\color{cCyan}\bfseries\fontsize{12}{14}\selectfont \faDatabase\ \ MAESTRÍA EN CIENCIA DE DATOS}
  };

  \node[anchor=center, text width=15cm, align=center] at ($(current page.center)+(0,1.2)$) {
    {\color{white}\bfseries\fontsize{40}{44}\selectfont Práctica 3}\\[0.35em]
    {\color{cMint}\bfseries\fontsize{20}{24}\selectfont \faSitemap\ Pipeline de Clasificación con Weka}
  };

  \draw[cCyan, line width=1.4pt] ($(current page.center)+(-3.6,-0.55)$) -- ++(7.2,0);

  \node[anchor=center, text width=13cm, align=center] at ($(current page.center)+(0,-1.3)$) {
    {\color{white!85!black}\fontsize{11}{14}\selectfont
     Construcción visual de un flujo de Machine Learning en KnowledgeFlow\\
     para clasificar el dataset Iris con cuatro algoritmos distintos}
  };

  \node[anchor=south, text width=14cm, align=center] at ($(current page.south)+(0,2.6)$) {
    \begin{tabular}{r l}
      \color{cCyan}\bfseries \faUser\ Elaborado por: & \color{white} Miguel Deza \\[3pt]
      \color{cCyan}\bfseries \faGraduationCap\ Programa: & \color{white} Maestría en Ciencia de Datos \\[3pt]
      \color{cCyan}\bfseries \faBook\ Curso: & \color{white} Introducción a Ciencia de Datos \\[3pt]
      \color{cCyan}\bfseries \faChalkboardTeacher\ Docente: & \color{white} Edgar Sarmiento Calisaya \\[3pt]
      \color{cCyan}\bfseries \faIcon{calendar-alt}\ Fecha: & \color{white} 05 de septiembre de 2026 \\
    \end{tabular}
  };
\end{tikzpicture}
\end{titlepage}
\restoregeometry

% -------------------- ÍNDICE --------------------
\tableofcontents
\thispagestyle{fancy}
\newpage

% -------------------- 1. INTRODUCCIÓN --------------------
\section{\faLightbulb\ \ Introducción}
Un proyecto de ciencia de datos no termina en limpiar datos: el siguiente paso suele ser construir un \textbf{modelo} que aprenda un patrón a partir de ejemplos ya etiquetados, para luego predecir la etiqueta de casos nuevos. A esto se le llama \textbf{clasificación}, uno de los problemas centrales del \textit{Machine Learning}.

En esta práctica utilicé \textbf{Weka} (\textit{Waikato Environment for Knowledge Analysis}), una herramienta de código abierto para minería de datos, y específicamente su entorno \textbf{KnowledgeFlow}: un editor visual donde en vez de escribir código se arrastran componentes y se conectan con flechas para armar un flujo de procesamiento, de forma parecida a como armé el flujo ETL de Power BI en la práctica anterior.

Trabajé con el dataset clásico \textbf{Iris} (Fisher, 1936): 150 flores de 3 especies distintas, descritas por 4 medidas numéricas. El objetivo fue construir un pipeline que cargue estos datos, los explore, entrene varios algoritmos de clasificación y compare qué tan bien predicen la especie correcta.

% -------------------- 2. OBJETIVO --------------------
\section{\faBullseye\ \ Objetivo}
\begin{objetivobox}
Entender las etapas de un proyecto de ciencia de datos a través de la construcción de un pipeline visual para clasificar un conjunto de datos usando Weka.
\end{objetivobox}

% -------------------- 3. HERRAMIENTAS --------------------
\section{\faToolbox\ \ Herramientas y fuentes de datos}
Para desarrollar esta práctica utilicé:
\begin{itemize}[leftmargin=1.6em, itemsep=5pt]
  \raggedright
  \item[\textcolor{cCyan}{\faDesktop}] \textbf{Weka 3.8.6} -- \url{https://www.cs.waikato.ac.nz/ml/weka} (instalado en Linux vía AUR)
  \item[\textcolor{cCyan}{\faSitemap}] \textbf{KnowledgeFlow Environment} -- editor visual de pipelines de Weka
  \item[\textcolor{cCyan}{\faTerminal}] \textbf{Simple CLI} -- línea de comandos de Weka, usada para el cuarto clasificador (ver Sección~\ref{sec:naivebayes})
  \item[\textcolor{cCyan}{\faDatabase}] \textbf{Dataset:} \texttt{iris.arff} -- 150 instancias, 4 atributos numéricos + clase (\texttt{Iris-setosa}, \texttt{Iris-versicolor}, \texttt{Iris-virginica})
\end{itemize}

% -------------------- 4. DESARROLLO --------------------
\section{\faCogs\ \ Desarrollo}

\subsection{\faSearch\ \ Exploración del dataset}
Antes de clasificar, analicé los 4 atributos numéricos del dataset (\texttt{sepallength}, \texttt{sepalwidth}, \texttt{petallength}, \texttt{petalwidth}, todos de tipo \textbf{REAL}) junto con la clase (\texttt{class}, tipo \textbf{nominal}, 3 categorías). El dataset tiene \textbf{150 instancias}, sin valores nulos, con la clase perfectamente balanceada (50 por especie). Sí detecté \textbf{3 filas exactamente duplicadas}, un detalle normal en este dataset tan estudiado (singularidad).

\begin{table}[H]
\centering
\small
\begin{tabular}{lrrrrr}
\toprule
\textbf{Atributo} & \textbf{Rango} & \textbf{Media} & \textbf{Var.} & \textbf{Desv. Est.} & \textbf{Correl. con clase} \\
\midrule
sepal length & 4.3 -- 7.9 & 5.84 & 0.68 & 0.83 & 0.78 \\
sepal width  & 2.0 -- 4.4 & 3.05 & 0.19 & 0.43 & -0.42 \\
petal length & 1.0 -- 6.9 & 3.76 & 3.09 & 1.76 & \textbf{0.95} \\
petal width  & 0.1 -- 2.5 & 1.20 & 0.58 & 0.76 & \textbf{0.96} \\
\bottomrule
\end{tabular}
\caption{Estadísticas descriptivas de los 4 atributos numéricos del dataset Iris.}
\end{table}

El siguiente gráfico deja ver con claridad que el tamaño del \textbf{pétalo} (largo y ancho) predice la especie mucho mejor que el tamaño del \textbf{sépalo} -- un dato que después confirmé al revisar el árbol de decisión del clasificador J48.

\figuraMarcada{img/grafico-correlacion-atributos.png}{0.95}{Correlación de cada atributo con la especie de iris. Los atributos del pétalo dominan claramente sobre los del sépalo.}

\textbf{Relaciones entre pares de atributos.} Más allá de qué tan bien predice cada atributo la especie, también es útil ver cómo se relacionan los atributos \emph{entre sí}. La matriz de correlación siguiente muestra que \texttt{petallength} y \texttt{petalwidth} están fuertemente correlacionados entre ellos (0.96) -- tiene sentido, ambos crecen juntos según el tamaño general de la flor -- y que \texttt{sepalwidth} es el único atributo que se correlaciona negativamente con todos los demás.

\figuraMarcada{img/grafico-correlacion-pares.png}{0.85}{Matriz de correlación entre los 4 atributos numéricos (relaciones entre pares).}

\textbf{Distribución de valores por especie.} Para entender mejor cómo se distribuyen los valores de los atributos clave (pétalo), separé sus valores por especie:

\figuraMarcada{img/grafico-distribucion-por-especie.png}{0.98}{Distribución (mínimo, cuartiles, mediana y máximo) del largo y ancho de pétalo, por especie.}

Estas cajas muestran visualmente lo mismo que ya sugería la correlación: \textit{Iris-setosa} ocupa un rango de valores completamente separado del resto (sin solaparse ni un poco), mientras que las cajas de \textit{versicolor} y \textit{virginica} sí se superponen parcialmente -- exactamente donde después aparecieron los errores de clasificación de los cuatro algoritmos.

\subsection{\faSitemap\ \ Construcción del pipeline visual}
En el KnowledgeFlow armé el siguiente flujo:

\begin{enumerate}[leftmargin=1.6em, itemsep=3pt, label=\textcolor{cTealDeep}{\textbf{\arabic*.}}]
  \item \textbf{ArffLoader} -- cargó el archivo \texttt{iris.arff}.
  \item \textbf{ClassAssigner} -- designó la columna \texttt{class} como la etiqueta a predecir.
  \item \textbf{CrossValidationFoldMaker} -- dividió los datos en 10 pliegues para entrenar y evaluar de forma honesta (nunca examinando al modelo con los mismos datos que usó para aprender).
  \item Cuatro ramas paralelas, una por cada \textbf{algoritmo de clasificación} (detallados en la siguiente sección), cada una con su propio \textbf{ClassifierPerformanceEvaluator} y \textbf{TextViewer} para revisar sus métricas por separado.
\end{enumerate}

\figuraMarcada{img/pipeline-completo.png}{0.98}{Pipeline completo armado en KnowledgeFlow: un origen de datos, validación cruzada, y cuatro algoritmos de clasificación evaluados en paralelo.}

\subsection{\faCogs\ \ Algoritmos de clasificación evaluados}
Ejecuté y comparé \textbf{cuatro} algoritmos con enfoques muy distintos entre sí:

\begin{itemize}[leftmargin=1.9em, itemsep=6pt]
  \raggedright
  \item[\textcolor{cTealDeep}{\faTree}] \textbf{J48:} construye un único árbol de decisión, con reglas del tipo ``si el pétalo mide menos de X, entonces es tal especie''.
  \item[\textcolor{cTealDeep}{\faLayerGroup}] \textbf{Random Forest:} entrena 100 árboles distintos sobre porciones aleatorias de los datos, y decide la especie por votación mayoritaria entre todos ellos.
  \item[\textcolor{cTealDeep}{\faBrain}] \textbf{Multilayer Perceptron:} una red neuronal que ajusta sus pesos internos durante 500 iteraciones para aprender fronteras de decisión más complejas.
  \item[\textcolor{cTealDeep}{\faDice}] \textbf{Naive Bayes:} un clasificador probabilístico -- calcula, para cada especie, qué tan probable es observar esas 4 medidas, y elige la especie más probable. Es el único de los cuatro que no aprende reglas ni pesos, sino una distribución estadística por atributo.
\end{itemize}

\figuraMarcada{img/resultado-random-forest.png}{0.95}{Resultado de Random Forest: 95.33\% de precisión (143/150 instancias correctas).}

\figuraMarcada{img/resultado-j48.png}{0.95}{Resultado de J48: 96\% de precisión (144/150 instancias correctas).}

\figuraMarcada{img/resultado-multilayer-perceptron.png}{0.95}{Resultado de Multilayer Perceptron: 97.33\% de precisión (146/150 instancias correctas), el mejor de los cuatro.}

\subsubsection{Cuarto clasificador: Naive Bayes}
\label{sec:naivebayes}
Agregué \textbf{NaiveBayes} (\texttt{Classifiers → bayes}) al mismo pipeline visual, conectado igual que los otros tres a la salida de \texttt{CrossValidationFoldMaker}. Antes de correrlo en KnowledgeFlow, verifiqué el algoritmo por la \textbf{Simple CLI} de Weka (\texttt{NaiveBayes -t iris.arff -x 10}, otra de las 4 aplicaciones oficiales de Weka) -- el resultado del pipeline visual coincidió exactamente con la CLI, confirmando que ambas vías evalúan con la misma metodología.

\figuraMarcada{img/resultado-naive-bayes.png}{0.95}{Resultado de Naive Bayes: 96\% de precisión (144/150 instancias correctas), empatado con J48.}

\subsection{\faChartBar\ \ Comparación de resultados}
Reuniendo las métricas de los cuatro algoritmos:

\begin{table}[H]
\centering
\small
\begin{tabular}{lrrrrr}
\toprule
\textbf{Algoritmo} & \textbf{Correctas} & \textbf{Precisión} & \textbf{Kappa} & \textbf{MAE} & \textbf{RMSE} \\
\midrule
Random Forest          & 143/150 & 95.33\,\% & 0.93 & 0.0408 & 0.1621 \\
J48                     & 144/150 & 96.00\,\% & 0.94 & 0.0350 & 0.1586 \\
Multilayer Perceptron   & \textbf{146/150} & \textbf{97.33\,\%} & \textbf{0.96} & \textbf{0.0327} & \textbf{0.1291} \\
Naive Bayes             & 144/150 & 96.00\,\% & 0.94 & 0.0342 & 0.1550 \\
\bottomrule
\end{tabular}
\caption{Tabla comparativa del rendimiento de los 4 algoritmos de clasificación (validación cruzada, 10 pliegues).}
\end{table}

\figuraMarcada{img/grafico-comparacion-precision.png}{0.95}{Precisión global de los cuatro algoritmos evaluados con validación cruzada de 10 pliegues.}

La diferencia entre los cuatro es pequeña (menos de 2 puntos porcentuales), pero la matriz de confusión de cada uno revela \emph{dónde} se equivoca cada algoritmo:

\figuraMarcada{img/grafico-errores-por-clase.png}{0.95}{Errores de clasificación por especie. Ningún algoritmo confunde jamás \texttt{Iris-setosa} con otra especie; todos los errores ocurren entre \texttt{versicolor} y \texttt{virginica}.}

Esto no es casualidad: el propio dataset documenta que \textit{setosa} es linealmente separable de las otras dos especies, mientras que \textit{versicolor} y \textit{virginica} no lo son entre sí. Los cuatro algoritmos ``descubrieron'' ese mismo patrón de forma independiente, sin importar su enfoque (árbol, ensamble, red neuronal o probabilístico).

\subsection{\faTree\ \ Visualización del árbol J48}
De los cuatro algoritmos, solo J48 se puede dibujar como un árbol simple (Random Forest son 100 árboles combinados, la red neuronal son pesos numéricos, y Naive Bayes es una tabla de probabilidades -- ninguno es un árbol único). Usé el componente \textbf{GraphViewer} conectado a la salida \texttt{graph} de J48 para visualizarlo:

\figuraMarcada{img/arbol-j48-graphviewer.png}{0.85}{Árbol de decisión generado por J48 en uno de los 10 pliegues de validación cruzada.}

El árbol confirma exactamente lo que ya había visto en la correlación de atributos: \textbf{nunca usa \texttt{sepallength} ni \texttt{sepalwidth}} para decidir -- todas sus preguntas son sobre \texttt{petalwidth} y \texttt{petallength}. La primera pregunta (¿\texttt{petalwidth} $\leq$ 0.6?) separa \textit{setosa} del resto sin ningún error; las preguntas siguientes van afinando la frontera entre \textit{versicolor} y \textit{virginica}.

% -------------------- 5. HALLAZGO TÉCNICO --------------------
\section{\faExclamationTriangle\ \ Hallazgo técnico}
\begin{incidenciabox}
Intenté conectar la salida \texttt{graph} también en Random Forest y en Multilayer Perceptron para visualizarlos igual que a J48, pero esa opción de conexión \textbf{no aparece} en el menú contextual de ninguno de los dos (solo ofrecen \texttt{batchClassifier}, \texttt{text} e \texttt{info}).

\textbf{Explicación:} la visualización tipo árbol solo está disponible para clasificadores que implementan la interfaz \texttt{Drawable} de Weka como un único árbol. Random Forest son 100 árboles combinados (no uno solo) y Multilayer Perceptron es una red de pesos numéricos -- ninguno de los dos se puede representar como un árbol de decisión simple, así que Weka correctamente no ofrece esa opción para ellos.
\end{incidenciabox}

% -------------------- 6. EVIDENCIA --------------------
\section{\faFileArchive\ \ Evidencia de ejecución}
\begin{evidenciabox}
Ejecuté el pipeline completo de principio a fin sin errores: los 16 componentes del flujo (carga, asignación de clase, validación cruzada, 4 clasificadores, 4 evaluadores, 4 visores de texto y el GraphViewer) terminaron correctamente.
\end{evidenciabox}

\figuraMarcada{img/evidencia-ejecucion-completa.png}{0.85}{Registro de estado (Status) de KnowledgeFlow: todos los componentes del pipeline terminaron con \texttt{Finished.}}

% -------------------- 7. CONCLUSIONES --------------------
\section{\faCheckCircle\ \ Conclusiones}
\begin{itemize}[leftmargin=1.6em, itemsep=6pt]
  \raggedright
  \item[\textcolor{cMint}{\faCheckCircle}] Construí un pipeline visual completo de ciencia de datos en KnowledgeFlow, desde la carga de datos hasta la comparación de múltiples modelos, sin escribir una sola línea de código.
  \item[\textcolor{cMint}{\faCheckCircle}] La validación cruzada de 10 pliegues me permitió comparar los algoritmos de forma honesta, evaluando cada uno con datos que nunca usó para entrenar.
  \item[\textcolor{cMint}{\faCheckCircle}] \textbf{Multilayer Perceptron} obtuvo la mejor precisión (97.33\%), seguido de un empate entre \textbf{J48} y \textbf{Naive Bayes} (96\% cada uno) y \textbf{Random Forest} (95.33\%) -- aunque la red neuronal fue notablemente más lenta de entrenar.
  \item[\textcolor{cMint}{\faCheckCircle}] Que J48 (un árbol de reglas) y Naive Bayes (probabilidades) empaten en precisión, pese a ser enfoques completamente distintos, muestra que para un dataset simple y bien separado como Iris, la elección del algoritmo importa menos que la calidad de los datos mismos.
  \item[\textcolor{cMint}{\faCheckCircle}] Las estadísticas descriptivas (correlación de atributos), el árbol de J48 y las matrices de confusión de los cuatro algoritmos contaron exactamente la misma historia de forma independiente: el tamaño del pétalo es lo que realmente distingue a las especies, y \textit{Iris-setosa} es trivialmente separable mientras que \textit{versicolor} y \textit{virginica} se solapan.
  \item[\textcolor{cMint}{\faCheckCircle}] Entendí una limitación real de las herramientas de visualización de modelos: no todo algoritmo se puede dibujar como árbol, y reconocer por qué (ensamble, red neuronal, probabilístico vs. árbol único) es tan importante como saber usar la herramienta.
  \item[\textcolor{cMint}{\faCheckCircle}] Aprendí que Weka no es solo KnowledgeFlow: usé la Simple CLI para verificar rápido un cuarto algoritmo con la misma metodología (validación cruzada de 10 pliegues) antes de integrarlo también al pipeline visual, confirmando que ambas vías dan el mismo resultado.
\end{itemize}

% -------------------- REFERENCIAS --------------------
\section{\faLink\ \ Referencias}
\begin{itemize}[leftmargin=1.6em, itemsep=4pt]
  \raggedright
  \item[\textcolor{cCyan}{\faLink}] Weka -- \url{https://www.cs.waikato.ac.nz/ml/weka}
  \item[\textcolor{cCyan}{\faLink}] Fisher, R.A. (1936). \textit{The use of multiple measurements in taxonomic problems} -- fuente original del dataset Iris.
\end{itemize}

\end{document}
